% s-253-p253-m003; original PDF 253, 254
\Aphorism{s-253-p253-m003}{350}
Da das Beste an die Jugend doch verschwendet wäre, so finde ich es recht und billig, wenn sie jene Mittelsorte geistiger Nahrung bevorzugt, die, ohne gerade verdorben zu sein, doch von keiner wählerischen Zunge mehr angerührt wird: das wohlfeil Pathetische, das unwissenschaftlich Empfindsame, das beschränkt Entrüstete. Schon in jungen Jahren kommt man oft in den Fall, Sachen nicht mehr lesen zu können, die man als Schüler hätte lesen sollen; es giebt nämlich, und in der ganz klassischen Litteratur, Dinge, für die man mit zwanzig Jahren schon zu alt ist. Um so wünschenswerther ist, dass man später öconomischer verfahre und das, was man immerhin einmal gelesen haben soll, in einem Alter lese, wo es Einem noch Vergnügen macht: so muss der Zwanzigjährige dem Dreissigjährigen, dieser dem Vierzigjährigen vorarbeiten und Geschmacksanachronismen ersparen.
\par

Diese Anordnung der Dinge wird von der Jugend übrigens instinctiv befolgt, während das begleitende Bewusstsein noch ganz im Schwärmerischen irregeht und die litterarischen \Foreign{liaisons} für genau so ewig und unzerstörbar hält wie der Jüngling seine erste Liebe. Ach! der Bücher und Kunstwerke, mit denen man sich auf Lebenszeit verheirathen kann, giebt es zum Glück nur wenige — und keineswegs sind es die leidenschaftlich bezaubernden, die, in welche man sich \Foreign{prima vista} verliebt. Was wissen wir heute von den Frauen und Büchern, die uns auf der Mittagshöhe des Lebens begegnen und den langen Abstieg begleiten werden? Aber es lohnt sich, auf sie zu warten.
\par

